\section{Introduction}

Generative recommendation replaces ranking over a catalog with sequence generation over item codes. This reformulation has become increasingly attractive because it allows recommender systems to inherit the training and decoding machinery of language models. Recent systems such as TIGER, LCRec, and MiniOneRec demonstrate that a discrete item tokenizer can be the bridge between item semantics and autoregressive recommendation \citep{rajput2023tiger,zheng2023lcrec,kong2025minionerec}. Once the tokenizer is fixed, however, its structural biases are inherited by every downstream stage, including constrained decoding, supervised fine-tuning, and recommendation-oriented reinforcement learning.

This makes tokenizer quality a first-order problem rather than a preprocessing detail. In our current MiniOneRec reproduction line, we repeatedly observed a pattern that does not fit the usual ``collision is the core problem'' narrative. After the official \texttt{sid-generate} collision-repair stage, the final code collision rate is already at the sub-percent level, yet many prediction errors remain. A closer look shows that these errors often occur after the model has already entered the correct semantic region. In other words, the prefix is frequently correct, but the final leaf token is not. This points to a more specific bottleneck: \emph{local ambiguity inside semantically plausible subtrees}.

At the same time, existing collaborative tokenizer methods suggest that collaborative signals should influence item tokenization, but they often do so through a single globally fused collaborative representation \citep{fang2026prism,liang2026resid,wang2026pit,liu2025discrec}. That design is natural for flat item representations, but it is less natural for hierarchical semantic IDs. Different levels in a semantic ID serve different structural roles. Upper levels act as coarse routing decisions, while deeper levels distinguish among increasingly similar candidates. A single collaborative view injected uniformly across all levels may therefore be both under-specified and unstable.

This paper develops the current state of that idea into a concrete tokenizer-side method. We present \textbf{MGR-SID}, a hierarchy-aware graph-regularized semantic ID framework that uses graph structure as the carrier of collaborative information. Rather than replacing the semantic tokenizer, MGR-SID keeps MiniOneRec's semantic \rqvae{} backbone and augments it with three train-only graph views: a purified coarse collaborative graph, a middle-resolution spectral graph, and a purified local transition graph. In the current \textbf{v1} implementation, we instantiate hierarchy-awareness with a simple but explicit level-wise rule: coarse regularization at level 1, spectral middle-view regularization at level 2, and local transition regularization at level 3.

Our evidence proceeds in three steps. First, we run diagnostic probes on Industrial and Office and show that middle-resolution graph views are indispensable: the best spectral middle view substantially outperforms heuristic middle views on same-$l_2$ ambiguity buckets. Second, we trace the provenance of the original MiniOneRec tokenizer and show that fair tokenizer comparison requires the upstream \rqvae{} training regime rather than the default local configuration. Third, under this aligned training setup, we show that hierarchy-aware graph regularization improves the \emph{final generated SID} even though it does not dominate raw train-stage collision.

The current draft is intentionally scoped to the tokenizer stage. We do not yet claim downstream SFT or RL gains. Instead, the goal is to establish a precise tokenizer-side story before carrying the new SID into the rest of the pipeline.

\paragraph{Contributions.} This draft makes three contributions.
\begin{enumerate}[leftmargin=1.4em]
  \item We identify a tokenizer-side bottleneck for MiniOneRec: after fair reproduction of the upstream pipeline, the remaining difficulty lies in local leaf ambiguity rather than in unrepaired global collision.
  \item We propose MGR-SID, a hierarchy-aware graph-regularized semantic ID framework that injects graph-structured collaboration into a semantic tokenizer through level-specific structural constraints instead of uniform feature fusion.
  \item We provide preliminary but concrete evidence for this design. On Industrial, under upstream-aligned \rqvae{} training, hierarchy-aware regularization reduces final SID collision from $0.00353$ to $0.00326$ and improves local same-$l_2$ ambiguity metrics, while uniform graph regularization is clearly harmful.
\end{enumerate}

The remainder of the paper is organized as follows. \Cref{sec:related} situates MGR-SID in prior work. \Cref{sec:background} summarizes the tokenizer pipeline and the empirical bottlenecks that motivate our design. \Cref{sec:method} presents the current MGR-SID v1 method. \Cref{sec:experiments} reports probe studies, provenance alignment, and tokenizer-stage results. \Cref{sec:conclusion} concludes with limitations and the next steps toward downstream SFT evaluation.
